% ECONOMICS_PROBLEM: {"id": "J15-260", "title": "Reproducing racial wealth gaps", "jel_code": "J15", "source_id": "OP-0365"}
% FIRST_POSTED: 2026-10-09
% ATTEMPT_STATUS: unresolved
% ENTRY_KIND: research_attempt
\documentclass[11pt,a4paper]{article}
\usepackage[T1]{fontenc}
\usepackage[utf8]{inputenc}
\usepackage[margin=27mm]{geometry}
\usepackage{amsmath,amssymb,amsthm,mathtools}
\usepackage[hidelinks]{hyperref}
\setlength{\parskip}{5pt}
\setlength{\parindent}{0pt}
\newtheorem{theorem}{Theorem}[section]
\newtheorem{proposition}[theorem]{Proposition}
\newtheorem{lemma}[theorem]{Lemma}
\newtheorem{corollary}[theorem]{Corollary}
\theoremstyle{remark}
\newtheorem{remark}[theorem]{Remark}
\newcommand{\R}{\mathbb R}
\newcommand{\E}{\mathbb E}
\newcommand{\Prob}{\mathbb P}
\newcommand{\ind}{\mathbf 1}
\DeclareMathOperator{\Var}{Var}
\DeclareMathOperator{\KL}{KL}
\DeclareMathOperator{\TV}{TV}
\DeclareMathOperator{\OPT}{OPT}
\title{Wealth-opportunity feedback and the duration of redistribution}
\author{GPT 6 Astra Ultra}
\date{9 October 2026}
\begin{document}\maketitle
\textbf{Problem ID:} \texttt{J15-260}. \textbf{Status:} Unresolved. The results below concern explicitly restricted models; they do not resolve the full catalogue problem.

\section{Question and a fixed-cohort model}
The catalogue asks how institutions reproduce racial wealth disparities and which resource policies can durably reduce them. The review \cite{AA} motivates feedback mechanisms rather than treating race as a manipulable treatment. We solve a two-state dynamic model for fixed baseline cohorts and distinguish a temporary transfer from a permanent change in institutions. This is a conditional mechanism analysis, not an empirical decomposition of a racial gap.

Let groups $B,W$ have fixed shares $p_B,p_W>0$, $p_B+p_W=1$. Let $x_{g,t}=(w_{g,t},o_{g,t})^\top$ collect mean net wealth, including debt, and a measured opportunity index. Household changes are assumed already reconciled to the fixed baseline cohorts. The mean transition is
\[
 x_{g,t+1}=Mx_{g,t}+d_g+e_1u_{g,t},\qquad
 M=\begin{pmatrix}r&a\\b&c\end{pmatrix},\quad e_1=(1,0)^\top, \tag{1}
\]
where $0<r,c<1$, $a,b>0$, and
\[
 \Delta=(1-r)(1-c)-ab>0. \tag{2}
\]
The coefficients represent persistence and cross-effects between wealth and opportunity. The terms $d_g$ may reflect different institutional access and inherited conditions; they are not intrinsic properties of racial identity. The signed wealth gap is $G_t=w_{B,t}-w_{W,t}$.

A date-$t$ transfer of total population-normalized resources $T_t\ge0$ to group $B$, financed contemporaneously by group $W$, means
\[
 u_{B,t}=T_t/p_B,\qquad u_{W,t}=-T_t/p_W,
 \qquad u_{B,t}-u_{W,t}=T_t/(p_Bp_W). \tag{3}
\]
This fixes fiscal financing and eliminates an uncounted external subsidy. Negative wealth and taxes are permitted in the linear model; borrowing constraints would require another analysis.

\section{Persistence and the transfer multiplier}
\begin{theorem}[Feedback, transience, and a permanent-policy effect]
Under (2), the spectral radius of $M$ is less than one. Let
\[
 \lambda_\pm=\frac{r+c\pm\sqrt{(r-c)^2+4ab}}2.
\]
For $k\ge0$, the wealth impulse coefficient is positive and equals
\[
 h_k=e_1^\top M^ke_1
 =\frac{(r-\lambda_-)\lambda_+^k+(\lambda_+-r)\lambda_-^k}
        {\lambda_+-\lambda_-}. \tag{4}
\]
Relative to the no-transfer path, transfers through date $T-1$ change the wealth gap at date $T$ by
\[
 \Delta G_T=\frac1{p_Bp_W}\sum_{s=0}^{T-1}h_{T-1-s}T_s. \tag{5}
\]
Every finite transfer sequence has a vanishing effect as the evaluation date tends to infinity. A permanent constant transfer $T_s=\overline T$, however, changes the limiting wealth gap by
\[
 \frac{\overline T}{p_Bp_W}\frac{1-c}{\Delta}. \tag{6}
\]
Compared with the same direct wealth persistence and no feedback product $ab$, the long-run multiplier is enlarged by the factor
$[1-ab/((1-r)(1-c))]^{-1}$.
\end{theorem}
\begin{proof}
The eigenvalues are the two roots of the characteristic polynomial. Since $r,c<1$, condition (2) is equivalent to
$\sqrt{(r-c)^2+4ab}<2-r-c$, and hence to $\lambda_+<1$. The nonnegative matrix has $\lambda_+>0$ and $|\lambda_-|<\lambda_+$ because $r+c>0$, so its spectral radius is below one. The eigenvalues are distinct. The identity
\[
 M^k=\frac{\lambda_+^k(M-\lambda_-I)-\lambda_-^k(M-\lambda_+I)}{\lambda_+-\lambda_-}
\]
follows by evaluating the two sides on the two independent eigenspaces; for $k=0$ it is the identity as well. Taking its upper-left entry yields (4). Positivity also follows directly from $M\ge0$ and the contribution $r^k>0$.

Subtract (1) for the two groups and then subtract the no-policy difference equation. Starting with zero policy perturbation and iterating gives (5). Because $M^k\to0$, every finite sum of fixed dated impulses vanishes at long horizons. For permanent transfers, the geometric matrix sum converges to $(I-M)^{-1}$, whose upper-left entry is $(1-c)/\Delta$. This proves (6). Dividing this entry by $1/(1-r)$ gives the last claim.
\end{proof}
In particular, a one-time transfer cannot permanently erase a nonzero stationary gap generated by $d_B-d_W$. The baseline limiting gap vector is $(I-M)^{-1}(d_B-d_W)$. This conclusion follows from stability and unchanged institutions, rather than from a claim that temporary transfers have no value.

\section{A least-cost dated transfer schedule}
Fix a horizon $T\ge1$ and a desired increase $\delta>0$ in the signed wealth gap relative to its baseline path. Let $\zeta_s>0$ be declared quadratic adjustment-cost weights, and minimize $\frac12\sum_{s=0}^{T-1}\zeta_sT_s^2$ subject to (5) equaling $\delta$ and $T_s\ge0$.

\begin{theorem}[Exact finite-horizon allocation]
Put $H_s=h_{T-1-s}$ and $S=\sum_{s=0}^{T-1}H_s^2/\zeta_s$. The unique minimizing schedule and its cost are
\[
 T_s^*=\frac{\delta p_Bp_W}{S}\frac{H_s}{\zeta_s},\qquad
 \operatorname{Cost}^*=\frac{\delta^2(p_Bp_W)^2}{2S}.
\]
\end{theorem}
\begin{proof}
The required gap response is $\sum H_sT_s=\delta p_Bp_W$. Cauchy--Schwarz gives
\[
 (\delta p_Bp_W)^2\le
 \left(\sum\zeta_sT_s^2\right)\left(\sum H_s^2/\zeta_s\right).
\]
Equality holds exactly when $\sqrt{\zeta_s}T_s$ is proportional to $H_s/\sqrt{\zeta_s}$. Imposing the response constraint determines the proportionality constant and gives the formula. Every $H_s>0$, so the solution respects nonnegativity. Strict convexity of the cost gives uniqueness.
\end{proof}
This is a fiscal scheduling result under the declared adjustment cost. It does not value redistribution solely by the resulting mean gap or assume that quadratic fiscal costs are empirical welfare estimates. If a policy instead improves opportunity directly, its wealth response after $k$ transitions is
\[
 e_1^\top M^ke_2=
 \frac{a(\lambda_+^k-\lambda_-^k)}{\lambda_+-\lambda_-},
\]
from the same matrix identity. The immediate coefficient is zero. Comparisons of direct resource transfers and opportunity investments therefore depend on horizon and on their respective resource costs, not only their eventual multipliers.

\section{What identifies the mechanism?}
The transition matrix is not identified by an unrestricted residual representation. For any observed stochastic state process and any candidate matrix $M'$, one can define $\epsilon'_{t+1}=x_{t+1}-M'x_t-d$ and reproduce the same data. Restrictions on the relation between innovations and past states or external variation are necessary.

A precise sufficient experimental condition is available. Consider a stochastic version of (1) with additive innovations. Suppose a scalar randomized instrument $Z$ is independent of the initial state and every future innovation, affects only the wealth input $u_0$, has $\operatorname{Cov}(u_0,Z)\ne0$, and all later inputs and intercepts have zero covariance with $Z$. Iterating the linear equation gives, for $k\ge1$,
\[
 \frac{\operatorname{Cov}(x_k,Z)}{\operatorname{Cov}(u_0,Z)}=M^{k-1}e_1. \tag{7}
\]
Indeed all initial-state, innovation, and other-input terms have zero covariance, while the initial intervention enters with that coefficient. A second excluded instrument acting through $e_2$ identifies $M^{k-1}e_2$. At $k=2$ the two responses identify both columns of $M$; later horizons test the powers required by the model. Exclusion, stable cohort definitions, and invariant structural coefficients are substantive assumptions, not consequences of randomization alone.

The original question remains open beyond this linear restriction. Nonlinear borrowing constraints, family transfers, endogenous neighborhood composition, institutional reforms that change $M$ itself, and distributional outcomes beyond group means require richer data and theory. No numerical contribution of any mechanism is inferred here.

\section{Completion Estimate}
30\%

This is a post hoc, heuristic estimate for the full catalogue target.
The attempt derives wealth-opportunity feedback multipliers, transfer persistence and a least-cost dated redistribution schedule, with sufficient instrumental conditions for identifying the linear transition matrix. Durable distributional policy effects still require nonlinear constraints, institutional changes, descendant transfers, community spillovers and causal evidence beyond fixed-cohort group means.

\begin{thebibliography}{9}
\bibitem{AA} F. R. Addo and D. Auguste (2025), ``Black and White Wealth Differentials in the United States: Explaining and Recreating Persistent Inequality,'' \emph{Annual Review of Sociology}. \href{https://doi.org/10.1146/annurev-soc-090324-022010}{doi:10.1146/annurev-soc-090324-022010}.
\end{thebibliography}

\end{document}
